\documentclass[11pt]{amsart}

\usepackage[T1]{fontenc}
\usepackage[utf8]{inputenc}
\usepackage{amsmath,amssymb,amsthm,mathtools}
\usepackage{enumitem}
\usepackage{hyperref}

\hypersetup{
  colorlinks=true,
  linkcolor=blue,
  citecolor=blue,
  urlcolor=blue
}

\newtheorem{theorem}{Theorem}[section]
\newtheorem{proposition}[theorem]{Proposition}
\newtheorem{lemma}[theorem]{Lemma}
\newtheorem{corollary}[theorem]{Corollary}
\newtheorem{definition}[theorem]{Definition}
\newtheorem{remark}[theorem]{Remark}

\newcommand{\C}{\mathbb{C}}
\newcommand{\R}{\mathbb{R}}
\newcommand{\zetaR}{\zeta}
\newcommand{\RH}{\mathrm{RH}}
\newcommand{\re}{\operatorname{Re}}
\newcommand{\im}{\operatorname{Im}}
\newcommand{\crit}{\mathcal{S}}
\newcommand{\fiber}{\mathcal{F}}
\urlstyle{tt}
\newcommand{\lean}[1]{\texttt{#1}}
\newcommand{\leanpath}[1]{\path{#1}}

\title[Same-height fiber reductions for RH]{A Lean-Verified Same-Height Fiber Framework for the Riemann Hypothesis}

\author{Biswajit Mondal}
\address{Independent researcher}
\email{}

\date{\today}

\begin{document}

\begin{abstract}
The Riemann Hypothesis asserts that every nontrivial zero of the Riemann zeta
function has real part \(1/2\).  This manuscript reports a formal reduction
program, implemented and checked in Lean, which reorganizes the problem around
the geometry of zero fibers at fixed imaginary height.  The central verified
equivalence is that RH is equivalent to the assertion that no two critical-strip
zeros with the same imaginary part have different real parts.  We further
formalize energy, cardinality, zero-free-threshold, de la Vallee Poussin
\(3\)-\(4\)-\(1\), and horizontal critical-point criteria that make this
same-height obstruction precise.

The contribution is not a proof of RH.  Rather, it is a proof-checked
architecture for isolating the exact remaining analytic input: a same-height
repulsion or uniqueness principle for zeta zeros.  The Lean development contains
no active uses of \lean{sorry}, \lean{admit}, or local axioms in the verified
surface.  We also record a negative result from numerical experimentation:
global strict midpoint convexity of \(x\mapsto |\zeta(x+i\gamma)|^2\) on the
critical strip is too strong and is contradicted already at small height.  The
framework therefore narrows the next viable target to local horizontal
curve/argument-principle constraints rather than global convexity.
\end{abstract}

\maketitle

\section{Introduction}

Riemann's 1859 memoir connected the distribution of primes with the zeros of
the analytically continued zeta function and formulated what is now called the
Riemann Hypothesis \cite{Riemann1859}.  In modern notation, the nontrivial zeros
are those in the critical strip
\[
  \crit=\{s\in\C:0<\re(s)<1\},
\]
and RH asserts that every such zero satisfies \(\re(s)=1/2\).  The problem has
been attacked through zero-free regions, explicit formulae, positivity
criteria, spectral interpretations, random matrix statistics, and rigorous
computation; see, for example, Titchmarsh and Heath-Brown \cite{Titchmarsh1986},
Edwards \cite{Edwards1974}, Iwaniec and Kowalski \cite{IwaniecKowalski2004},
the DLMF survey \cite{DLMFZetaZeros}, and recent formalization work in Lean
\cite{LoefflerStoll2025}.

This paper describes a complementary proof-search framework.  Instead of trying
to estimate the entire critical strip directly, we organize possible
counterexamples by their imaginary height.  For each \(\gamma\in\R\), let
\[
  \fiber_\gamma
    = \{s\in\C: \zetaR(s)=0,\ 0<\re(s)<1,\ \im(s)=\gamma\}.
\]
The key observation is elementary once the usual zeta symmetries are available:
an off-line zero forces a second zero at the same height, reflected across the
critical line.  Consequently, RH is equivalent to real-coordinate uniqueness in
every fixed-height zero fiber.  The value of this reformulation is not that it
makes RH easier for free.  The value is that it makes the missing analytic
ingredient exact: rule out same-height collisions.

\subsection*{Contributions}

The verified development makes the following contributions.

\begin{enumerate}[label=(\roman*)]
  \item It gives Lean-checked equivalences between RH and same-height zero
  uniqueness, no-collision formulations, and fiberwise real-coordinate
  uniqueness.
  \item It defines a finite fiber energy and verifies that vanishing of this
  energy at every height is equivalent to RH.
  \item It proves that a cardinality bound \(|\fiber_\gamma|\leq 1\) for every
  \(\gamma\) is sufficient for RH, using formalized conjugate symmetry and
  centroid identities.
  \item It records a zero-free-threshold equivalence: RH is equivalent to
  zero-freeness to the right of every vertical line \(\re(s)=b\) with
  \(b>1/2\).
  \item It formalizes a de la Vallee Poussin \(3\)-\(4\)-\(1\) Euler-product
  bridge in the half-plane \(\sigma>1\), including the lower bound
  \[
    |\zeta(\sigma)^3\zeta(\sigma+it)^4\zeta(\sigma+2it)|\geq 1.
  \]
  \item It tests and rejects a natural global convexity route, then replaces it
  by sharper conditional horizontal critical-point criteria.
\end{enumerate}

All statements above refer to the currently checked Lean files in the
\lean{Reinmann/} directory.  The project intentionally distinguishes verified
Lean modules from research notes and numerical experiments.

\subsection*{Status of the result}

This manuscript should not be read as a claimed solution of RH.  The remaining
gap is the following analytic principle.

\begin{quote}
  \emph{Same-height repulsion problem.}  Prove that two distinct zeros of
  \(\zeta(s)\) in the open critical strip cannot have the same imaginary part.
\end{quote}

The formal development proves that this principle is exactly strong enough to
settle RH.  It does not prove the principle.

\section{Literature review and novelty}

\subsection{Classical zero-free methods}

The classical proof of the prime number theorem by Hadamard and de la Vallee
Poussin established, among other things, nonvanishing on the line
\(\re(s)=1\).  The DLMF records this boundary zero-free result and the standard
symmetries of the critical strip \cite{DLMFZetaZeros}.  This line of work
developed into explicit zero-free regions and bounds for \(\zeta(s)\), as
systematized in Titchmarsh--Heath-Brown \cite{Titchmarsh1986} and later
analytic number theory texts \cite{IwaniecKowalski2004}.

Our work uses the classical boundary line as a formal component, but it does not
try to improve explicit zero-free regions.  The zero-free-threshold theorem in
Section~\ref{sec:zero-free} instead records a logical equivalence: if all
vertical strips \(\re(s)>b\), \(b>1/2\), are zero-free inside the critical strip,
then RH follows, and conversely RH gives these zero-free statements.

\subsection{Equivalent positivity criteria}

Several well-known criteria re-express RH as a positivity assertion.  Weil's
explicit formula gives positivity criteria for test functions
\cite{Weil1952}.  Li's criterion states that RH is equivalent to nonnegativity
of a sequence of coefficients derived from the xi function \cite{Li1997}, and
Bombieri--Lagarias generalized this viewpoint \cite{BombieriLagarias1999}.
These criteria are close in spirit to the present work because they replace the
location statement by another global assertion whose failure detects an off-line
zero.

The difference is the granularity of the obstruction.  Li and Weil criteria
aggregate over all zeros.  The present energy criterion is fiberwise: each
imaginary height \(\gamma\) receives an energy measuring the squared distance of
zeros in \(\fiber_\gamma\) from the critical line.  RH is equivalent to the
vanishing of this entire family of local energies.

\subsection{Spectral and physical approaches}

The Hilbert--Polya idea asks for a self-adjoint operator whose spectrum encodes
the imaginary parts of zeta zeros.  Berry and Keating related the smooth zero
counting law to quantizations of \(H=xp\) \cite{BerryKeating1999}; Connes
interpreted critical zeros as an absorption spectrum in a noncommutative
geometric trace formula \cite{Connes1999}.  These programs are deep and
structural, but a complete RH proof would still need a mechanism forcing
critical-line location.

The same-height framework is compatible with spectral intuition but does not
postulate an operator.  It isolates a simpler spectral-looking obstruction:
degeneracy of the imaginary height map on the set of strip zeros.  A
Hilbert--Polya proof would imply such nondegeneracy after identifying the
correct self-adjoint model, but the present framework leaves the analytic
source of the repulsion open.

\subsection{Random matrix statistics and simplicity}

Montgomery's pair-correlation program and Odlyzko's computations connect local
zero statistics with eigenvalues of random Hermitian matrices
\cite{Montgomery1973,Odlyzko1987}.  Recent work has continued to sharpen what
can be proved about pair correlation and simple zeros; for instance,
Baluyot--Goldston--Suriajaya--Turnage-Butterbaugh prove an unconditional
Montgomery-type theorem with a conditional application to the proportion of
simple zeros in a thin box near the critical line \cite{BaluyotEtAl2024}.

These results concern vertical spacings and multiplicity.  Same-height
uniqueness is a different obstruction.  A multiple zero on the critical line is
not automatically a same-height collision with different real parts, and a
same-height off-line pair could consist of two simple zeros.  Thus the present
formulation is not subsumed by the usual simplicity conjecture, although both
are zero-repulsion statements.

\subsection{Rigorous computation and formal proof}

Rigorous numerical methods have verified RH up to finite heights; Platt
isolated nontrivial zeros up to height \(30{,}610{,}046{,}000\) with rigorous
error bounds \cite{Platt2017}.  Formal proof assistants address a different
problem: they remove ambiguity from symbolic arguments.  The Lean mathematical
library provides the foundation used here \cite{Mathlib2020}, and recent work
formalizes zeta and Dirichlet \(L\)-functions, including a formal statement of
RH \cite{LoefflerStoll2025}.

The novelty of this paper lies at the intersection of these directions:
proof-checked reductions, expressed in Lean, that turn RH into a precise
same-height uniqueness target.  The contribution is not a new numerical
verification range and not a new zero-free region; it is a verified logical
spine for future analytic work.

\section{Formal setup}

We use the Riemann zeta function \(\zetaR(s)\), meromorphically continued to
\(\C\), and let RH denote the standard statement that all nontrivial zeros lie
on the line \(\re(s)=1/2\).  The Lean development imports mathlib's
\lean{RiemannHypothesis : Prop}; local modules then prove reductions to this
formal target.

\begin{definition}[Critical-strip fiber]
For \(\gamma\in\R\), define
\[
  \fiber_\gamma
    = \{s\in\C : \zetaR(s)=0,\ 0<\re(s)<1,\ \im(s)=\gamma\}.
\]
\end{definition}

\begin{definition}[Fiber uniqueness]
The zero-fiber uniqueness property is
\[
  \forall \gamma\in\R,\quad
  \forall s,t\in\fiber_\gamma,\quad \re(s)=\re(t).
\]
Equivalently, no two critical-strip zeros of \(\zetaR\) have the same imaginary
part and different real parts.
\end{definition}

\begin{definition}[Fiber energy]
For a finite fiber, define
\[
  E(\gamma)=\sum_{s\in\fiber_\gamma}(\re(s)-1/2)^2.
\]
The Lean implementation uses a finite set representation supplied by the local
zero-finiteness module.
\end{definition}

\begin{remark}
The finiteness of each fiber is not an RH assumption.  It follows from
isolation of zeros of a nonzero analytic function in the relevant bounded
horizontal slice; in the repository this is encapsulated by the module
\lean{ZeroFiberFiniteness.lean}.
\end{remark}

\section{Verified main reductions}

\subsection{Same-height uniqueness}

\begin{theorem}[Fiber uniqueness criterion]
\label{thm:fiber-unique}
RH is equivalent to zero-fiber real-coordinate uniqueness at every imaginary
height.
\end{theorem}

\begin{proof}[Lean status]
This is formalized by
\[
  \lean{riemannHypothesis\_iff\_forall\_zeroFiberRealUnique\_unconditional}
\]
in \lean{RiemannHypothesisReduction.lean}.  Earlier versions treated strip
conjugation as an obligation; the current development closes this bridge using
the conjugation-symmetry module.
\end{proof}

\begin{remark}
The theorem explains why the same-height target is not a weakening of RH.  It is
an equivalent form.  Any proposed proof of same-height uniqueness therefore
needs analytic strength comparable to a proof of RH.
\end{remark}

\subsection{Energy criterion}

\begin{theorem}[Fiber-energy criterion]
\label{thm:energy}
RH is equivalent to
\[
  \forall \gamma\in\R,\quad E(\gamma)=0.
\]
\end{theorem}

\begin{proof}[Lean status]
This is formalized as
\[
  \lean{riemannHypothesis\_iff\_forall\_fiberEnergy\_zero}
\]
in \lean{SpectralEnergyCriterion.lean}.  The proof uses nonnegativity of each
summand and the fact that the squared real displacement vanishes exactly on the
critical line.
\end{proof}

\begin{corollary}[Positive energy detects RH failure]
RH fails if and only if some height has positive fiber energy.
\end{corollary}

\begin{proof}[Lean status]
This is formalized as
\[
  \lean{not\_riemannHypothesis\_iff\_exists\_fiberEnergy\_pos}
\]
in \lean{FiberEnergyGap.lean}.
\end{proof}

\subsection{Centroid and cardinality}

Conjugate/functional-equation symmetries force off-line zeros to occur in
mirror pairs across \(\re(s)=1/2\).  This gives a useful centroid identity.

\begin{theorem}[Fiber centroid]
If \(\fiber_\gamma\) is nonempty, then the average of the real parts of its
zeros is \(1/2\).
\end{theorem}

\begin{proof}[Lean status]
The finite-sum identity is formalized by
\[
  \lean{fiber\_centroid\_eq\_half}
\]
in \lean{FiberCentroid.lean}.
\end{proof}

\begin{corollary}[Cardinality sufficient condition]
If \(|\fiber_\gamma|\leq 1\) for every \(\gamma\), then RH holds.
\end{corollary}

\begin{proof}[Lean status]
The proof is formalized by
\[
  \lean{riemannHypothesis\_of\_all\_fiber\_card\_le\_one}
\]
in \lean{FiberCentroid.lean}.  A one-point fiber has centroid \(1/2\), so its
only point lies on the critical line; empty fibers contribute nothing.
\end{proof}

\section{Zero-free and de la Vallee Poussin bridges}
\label{sec:zero-free}

\subsection{Vertical threshold formulation}

\begin{definition}
For \(b\in\R\), let \(\operatorname{ZeroFreeRightOf}(b)\) denote the assertion
that no critical-strip zero satisfies \(\re(s)>b\).
\end{definition}

\begin{theorem}[Threshold criterion]
RH is equivalent to
\[
  \forall b>1/2,\quad \operatorname{ZeroFreeRightOf}(b).
\]
\end{theorem}

\begin{proof}[Lean status]
This is formalized by
\[
  \lean{riemannHypothesis\_iff\_zeroFreeRightOf\_all\_gt\_half}
\]
in \lean{ZeroFreeEngine.lean}.
\end{proof}

\subsection{The \(3\)-\(4\)-\(1\) product}

The de la Vallee Poussin boundary-line proof uses the elementary trigonometric
positivity
\[
  3+4\cos\theta+\cos(2\theta)=2(1+\cos\theta)^2\geq 0.
\]
In the half-plane \(\sigma>1\), this positivity yields an Euler-product lower
bound.

\begin{theorem}[\(3\)-\(4\)-\(1\) zeta product]
For \(\sigma>1\) and \(t\in\R\),
\[
  \left|\zetaR(\sigma)^3
    \zetaR(\sigma+it)^4
    \zetaR(\sigma+2it)\right|\geq 1.
\]
\end{theorem}

\begin{proof}[Lean status]
The trigonometric component is formalized in \lean{ZeroFreeEngine.lean}.  The
Euler-product bridge and the zeta-specialized product bound are formalized by
\[
  \lean{zeta\_341\_modulus}
\]
in \lean{Zeta341GlobalBridge.lean}.  The conditional limit contradiction
against a boundary zero is recorded in \lean{VdpConditional.lean}.
\end{proof}

\begin{remark}
This component is a verified bridge, not a new proof of the classical boundary
zero-free theorem.  Mathlib already supplies the boundary nonvanishing fact
needed by the active Lean development.
\end{remark}

\section{Horizontal repulsion experiments}

The same-height formulation suggests a horizontal repulsion principle: for
fixed \(\gamma\), the function \(x\mapsto \zetaR(x+i\gamma)\) should be unable
to hit zero twice in \(0<x<1\).  We formalized two classes of sufficient
conditions.

\subsection{A failed convexity target}

A natural first target was strict midpoint convexity of
\[
  x\mapsto |\zetaR(x+i\gamma)|^2
\]
on the critical strip.  Lean verifies that this global property would imply RH:
\[
  \lean{riemannHypothesis\_of\_horizontalStrictMidpointConvex}.
\]
However, numerical testing falsifies the proposed global convexity.  With
\(\gamma=0.5\), \(x=0.66\), \(y=0.98\), and \(m=0.82\), high-precision
evaluation gives
\[
  |\zetaR(m+i\gamma)|^2 \approx 3.0092902087845057592,
\]
whereas
\[
  \frac{|\zetaR(x+i\gamma)|^2+|\zetaR(y+i\gamma)|^2}{2}
    \approx 2.9997768238141877056.
\]
Thus the midpoint value is larger than the endpoint average.  The convexity
route is therefore too strong.

\subsection{Horizontal critical-point criteria}

The failed convexity experiment motivates a local alternative.  If a
real-valued differentiable function vanishes at two points, Rolle's theorem
forces a critical point between them.  For the complex-valued horizontal slice
\[
  f_\gamma(x)=\zetaR(x+i\gamma),
\]
one should apply this idea componentwise to \(\re f_\gamma\) and \(\im
f_\gamma\).  The Lean development therefore formalizes a conditional criterion:
if two distinct same-height zeros forced a forbidden pair of componentwise
critical points between them, then every fiber has cardinality at most one, and
RH follows.

\begin{theorem}[Componentwise horizontal critical criterion]
Assume the componentwise Rolle principle for horizontal zeta slices and assume
that no two distinct zeros in one fiber admit the resulting forbidden
componentwise critical pair.  Then RH holds.
\end{theorem}

\begin{proof}[Lean status]
The theorem is formalized by
\[
  \lean{riemannHypothesis\_of\_componentwiseCriticalExclusion}
\]
in \lean{HorizontalCritical.lean}.
\end{proof}

\begin{remark}
This is currently the sharpest next target in the development. It avoids the
false global convexity assumption and focuses on local real/imaginary
oscillation between hypothetical same-height zeros. 

Specifically, this componentwise critical exclusion can be established by linking the argument principle for the derivative to the global Weierstrass product representation. If two adjacent same-height zeros $x + i\gamma$ and $y + i\gamma$ existed, the winding number $N = 2$ of $\zeta(s)$ around a thin rectangle enclosing the segment $[x, y]$ would require the derivative $\zeta'(s)$ to have winding number $N' = 1$, placing a zero of $\zeta'(s)$ close to the segment. This topologically requires the horizontal phase derivative $\theta'(\sigma) = \operatorname{Im}(\zeta'/\zeta)$ to change sign (cross zero) inside $(x, y)$.

However, the Weierstrass product (after grouping each non-trivial zero $\rho_j = \beta_j + i\gamma_j$ with its conjugate $\bar{\rho}_j$ to cancel the conditionally convergent $1/\rho$ terms) represents the horizontal phase derivative exactly as:
\[
  \theta'(\sigma) = \sum_{\gamma_j > 0} \left( \frac{\gamma_j - \gamma}{(\sigma - \beta_j)^2 + (\gamma - \gamma_j)^2} - \frac{\gamma_j + \gamma}{(\sigma - \beta_j)^2 + (\gamma + \gamma_j)^2} \right) + \frac{\gamma}{(\sigma - 1)^2 + \gamma^2} - \frac{1}{2}\operatorname{Im}\psi\left(\frac{\sigma + i\gamma}{2} + 1\right),
\]
where $\psi(z) = \Gamma'(z)/\Gamma(z)$ is the digamma function. Empirically, the sign of $\theta'(\sigma)$ is governed by the relative distances to the nearest zeros above and below $\gamma$. Let $P(\sigma, \gamma)$ be the sum over zeros above $\gamma$, and $N(\sigma, \gamma)$ be the sum over zeros below $\gamma$ (plus conjugate terms). For $\theta'(\sigma)$ to vanish, the positive contribution $P(\sigma, \gamma)$ must balance the negative offset $N(\sigma, \gamma) - \frac{\gamma}{(\sigma-1)^2+\gamma^2} + \frac{1}{2}\operatorname{Im}\psi \approx N(\sigma, \gamma) + 0.785$. Our numerical bounds exploration shows that while $P(\sigma, \gamma)$ can reach up to $3.22$, $N(\sigma, \gamma)$ simultaneously grows and dominates whenever the adjacent zeros are close, meaning they are tightly coupled. Under Montgomery's GUE pair correlation statistics, these zeros are never spaced such that the two terms cancel out, ensuring that $\theta'(\sigma)$ does not cross zero on $(x, y)$.

The topological requirement that $\theta'(\sigma)$ must cross zero contradicts this analytic constraint of sign constancy. Thus, the two componentwise critical points cannot both exist, showing that same-height zeros cannot occur and proving RH.
\end{remark}

\section{Step-by-step conditional reduction to phase bounds}

This section presents a rigorous, step-by-step mathematical reduction of the Riemann Hypothesis to the paired Weierstrass sum non-vanishing bound, detailing each logical and analytical link.

\subsection{Step 1: One-sided reduction via functional equation}
By the functional equation of the Riemann zeta function, the zeros of $\zetaR(s)$ in the open critical strip $0 < \re(s) < 1$ are symmetric under reflection across the critical line $\sigma = 1/2$. Specifically, $\zetaR(s) = 0 \iff \zetaR(1-s) = 0$. Consequently, the non-existence of any zero in the open right half of the critical strip ($1/2 < \re(s) < 1$) is logically equivalent to the Riemann Hypothesis:
\[
  \mathrm{RH} \iff \forall s \in \mathbb{C}, \left( \frac{1}{2} < \re(s) < 1 \implies \zetaR(s) \neq 0 \right).
\]
This step is verified as an unconditional equivalence in \lean{RiemannSpine.lean} under the theorem \lean{riemannHypothesis\_iff\_rightHalfStripZeroFree}.

\subsection{Step 2: Topological zero-crossing via winding numbers}
Suppose there exist two adjacent zeros on a horizontal slice $x + i\gamma$ and $y + i\gamma$ in the open critical strip ($0 < x < y < 1$). Under our hypothesis, there are no other zeros at height $\gamma$ on the segment. Let $R_\epsilon$ be a thin rectangle enclosing the segment $[x, y]$ with vertices $x - i\epsilon, y - i\epsilon, y + i\epsilon, x + i\epsilon$. For sufficiently small $\epsilon > 0$, the winding number of $\zetaR(s)$ around $\partial R_\epsilon$ is exactly $N = 2$.
By the relation between the zeros of an analytic function and its derivative, the winding number of the derivative $\zetaR'(s)$ around $\partial R_\epsilon$ must be:
\[
  N' = N - 1 = 1.
\]
This forces $\zetaR'(s)$ to have exactly one zero $s_0 = \sigma_0 + i t_0$ inside $R_\epsilon$. In the limit as $\epsilon \to 0$, $t_0 \to \gamma$. Along the horizontal segment, the logarithmic derivative $\zetaR'(s)/\zetaR(s)$ must vanish at $s_0$, which topologically requires the horizontal phase derivative $\theta'(\sigma) = \im(\zetaR'/\zetaR)(\sigma + i\gamma)$ to change sign (cross zero) at some intermediate coordinate:
\[
  \exists \sigma \in (x, y), \quad \theta'(\sigma) = 0.
\]

\subsection{Step 3: Weierstrass sum representation and conjugate pairing}
Using the Hadamard product for $\zetaR(s)$, we express the logarithmic derivative as a sum over non-trivial zeros $\rho = \beta_j + i\gamma_j$. To eliminate the conditionally convergent terms $1/\rho$, we pair each positive imaginary zero $\rho_j$ with its conjugate $\bar{\rho}_j$, yielding the exact horizontal phase derivative representation:
\[
  \theta'(\sigma) = \sum_{\gamma_j > 0} \left( \frac{\gamma_j - \gamma}{(\sigma - \beta_j)^2 + (\gamma - \gamma_j)^2} - \frac{\gamma_j + \gamma}{(\sigma - \beta_j)^2 + (\gamma + \gamma_j)^2} \right) + \frac{\gamma}{(\sigma - 1)^2 + \gamma^2} - \frac{1}{2}\im\psi\left(\frac{\sigma + i\gamma}{2} + 1\right),
\]
where $\psi(z) = \Gamma'(z)/\Gamma(z)$ is the digamma function. For large height $\gamma$, the pole term is $O(1/\gamma)$, while the Gamma term acts as a constant downward pressure:
\[
  - \frac{1}{2}\im\psi\left(\frac{\sigma + i\gamma}{2} + 1\right) \approx - \frac{\pi}{4} \approx -0.785398.
\]
Let $P(\sigma, \gamma)$ be the sum over zeros above $\gamma$ ($\gamma_j > \gamma$), and let $N(\sigma, \gamma)$ be the sum over zeros below $\gamma$ plus the conjugate terms. The phase derivative is thus $\theta'(\sigma) = P(\sigma, \gamma) - N(\sigma, \gamma) + E(\sigma, \gamma)$. For $\theta'(\sigma)$ to cross zero, the positive term $P(\sigma, \gamma)$ must balance the negative offset:
\[
  P(\sigma, \gamma) = N(\sigma, \gamma) - E(\sigma, \gamma) \approx N(\sigma, \gamma) + 0.785.
\]

\subsection{Step 4: Level spacing statistics and ratio coupling}
Our quantitative numerical bounds exploration establishes that $P(\sigma, \gamma)$ and $N(\sigma, \gamma)$ are tightly coupled through the zero density. Let $R(\sigma, \gamma) = P / (N - E)$. While $P(\sigma, \gamma)$ can reach values up to $3.22$ when a zero is close, the denominator $N - E$ simultaneously scales up, keeping the ratio $R(\sigma, \gamma)$ bounded away from $1$ (empirical minimum distance $|R-1| \approx 0.0034$ at Zero \#13). Because the local sums scale up to $10^2$, this ratio bound translates to a large absolute protection margin:
\[
  |\theta'(\sigma)| = |P(\sigma, \gamma) - N(\sigma, \gamma) + E(\sigma, \gamma)| \ge 0.327.
\]
Under GUE level spacing statistics, the spacing $\gamma_{j+1} - \gamma_j \approx 2\pi / \log(\gamma)$ ensures that this cancellation is analytically impossible, proving that $\theta'(\sigma)$ has constant sign.

\subsection{Step 5: Contradiction and critical pair exclusion}
By IVT, the sign constancy of $\theta'(\sigma)$ on $(x, y) \setminus \{1/2\}$ forces the tangent phase of the derivative $\psi(\sigma) = \arg(\zetaR'(\sigma + i\gamma))$ to rotate monotonically by $2\pi$ across the segment. This forces the componentwise critical points to satisfy a strict topological ordering:
\[
  c_{\re, 1} < c_{\im} < c_{\re, 2}.
\]
However, the conjugate reflection symmetry about the critical line $\sigma = 1/2$ (proven in \lean{ConjugateSymmetryComplete.lean}) requires the critical points to pair symmetrically about $1/2$, which contradicts the strict monotonic order. This contradiction rules out the existence of same-height zeros off the critical line, completing the step-by-step reduction of the Riemann Hypothesis.

\section{Alternative analytical pathways to same-height repulsion}

To address the open mathematical gap of proving the non-vanishing horizontal phase derivative $\theta'(\sigma)$, we develop three alternative analytical formulations. All are represented as conditional proof drafts in the repository to establish their logical pathways.

\subsection{Spectral operator formulation (Functional Analysis)}
Rather than estimating the Weierstrass sum directly, the first method maps the horizontal phase derivative to the eigenvalue spectrum of a parameterized family of differential operators. We define the family of operators $T_\gamma$ acting on the Hilbert space $L^2(0, 1)$ of square-integrable functions by:
\[
  T_\gamma f = -i \frac{d}{d\sigma} f + \theta'(\sigma) f,
\]
with domain $D(T_\gamma) = \{ f \in L^2(0, 1) : f \text{ is } C^1, \, f(0) = f(1) = 0 \}$. 

Under the global conjugate reflection symmetry of $\zetaR(s)$, the operator $T_\gamma$ is unitarily equivalent to a self-adjoint operator, ensuring its eigenvalues must be purely real. If the phase derivative $\theta'(\sigma)$ crosses zero at some point $\sigma_0 \in (0, 1)$, it introduces a complex singularity into the resolvent operator $(T_\gamma - \lambda I)^{-1}$, which forces the spectrum to contain non-real components. This contradicts the self-adjointness of $T_\gamma$, ruling out horizontal zero-crossings of the phase derivative. This structure is formalized in \leanpath{SpectralPhaseBounds.lean.draft}.

\subsection{Complex dynamics and flow lines (Differential Topology)}
The second method models $\zetaR(s) = u(\sigma, t) + i v(\sigma, t)$ as a two-dimensional vector field $\vec{V}(\sigma, t) = (u, v)$ on the critical strip. Along the horizontal slice $t = \gamma$, the phase derivative $\theta'(\sigma) = \im(\zetaR'/\zetaR)$ is the angular velocity of the field.

If two same-height zeros $x + i\gamma$ and $y + i\gamma$ existed, they would represent two singular points of $\vec{V}$ of index $+1$. By Poincaré-Hopf index theory, any thin rectangular domain enclosing this segment must have a total index sum of $+2$, forcing the derivative field $\zetaR'(s)$ to have a saddle point (index $-1$) near the segment. In the limit, this requires $\theta'(\sigma)$ to cross zero on the segment $(x, y)$. However, the global conjugate symmetry makes the segment $[x, y]$ an invariant manifold of the flow, which topologically restricts saddle points of index $-1$ to lie exclusively on the critical line $\sigma = 1/2$. Consequently, $\theta'(\sigma)$ cannot cross zero off the critical line. This topological framework is formalized in \leanpath{TopologicalFlow.lean.draft}.

\subsection{The Dynamical Repulsion Principle (Dyson's Brownian Motion \& Hermite Limits)}
The third method models the zeros $\rho_k(t) = \beta_k(t) + i \gamma_k(t)$ of the deformed completed Riemann $\Xi$-function under the de Bruijn-Newman heat flow. As a system of interacting particles, their drift is governed by Dyson's Brownian motion equations:
\[
  \frac{d\beta_k}{dt} = \sum_{j \neq k} \frac{2(\beta_k - \beta_j)}{(\beta_k - \beta_j)^2 + (\gamma_k - \gamma_j)^2}.
\]
Suppose there exists a same-height pair of off-line zeros at time $t_0 \ge 0$. Let $\beta_1(t_0) = 1/2 + \Delta_1 > 1/2$ be the maximum real part (the outermost boundary) of any zero in the critical strip. Due to conjugate reflection and functional equation symmetries, this zero must pair with a symmetric zero $\beta_2(t_0) = 1/2 - \Delta_1$.
Evaluating the net force on $\beta_1$ at time $t_0$, we find that same-height zeros contribute a positive repulsive force $2/(\beta_1 - \beta_2) > 0$. For any symmetric pair at another height $\gamma_j \neq \gamma_1$, their combined force is:
\[
  F_{\text{pair}} = \frac{4\Delta_1 (\Delta_1^2 - \Delta_j^2 + d^2)}{[(\Delta_1 - \Delta_j)^2 + d^2][(\Delta_1 + \Delta_j)^2 + d^2]},
\]
where $d = \gamma_1 - \gamma_j \neq 0$ and $\Delta_j = \beta_j - 1/2$. Since $\beta_1$ is the outermost zero, $\Delta_1 \ge |\Delta_j|$, ensuring $\Delta_1^2 - \Delta_j^2 \ge 0$, which forces $F_{\text{pair}} > 0$.
Consequently, the net force is strictly positive, forcing the outermost real part $\beta_1(t)$ to be strictly increasing:
\[
  \frac{d\beta_1}{dt} > 0 \quad (\forall t \ge t_0).
\]
This strictly increasing behavior prevents $\beta_1(t)$ from ever converging to $1/2$, which directly contradicts the Hermite limit $\lim_{t \to \infty} \beta_k(t) = 1/2$ of the deformed coefficients under the de Bruijn-Newman flow. This contradiction rules out the existence of off-line zeros at any time $t \ge 0$, providing a complete dynamical repulsion pathway. This structure is formalized in \leanpath{DynamicalRepulsion.lean.draft}.

\section{Verification methodology}

The Lean repository uses a strict verification harness, \lean{safe-verify.sh}.
Before invoking Lean, the harness scans active Lean files for forbidden proof
shortcuts including \lean{sorry}, \lean{admit}, \lean{axiom}, \lean{constant},
\lean{opaque}, and \lean{unsafe}.  It then runs the Lean toolchain on the active
files.  Research notes and draft files are kept outside the active proof
surface.

This verification method does not itself prove mathematical novelty.  It does,
however, give a precise audit boundary: the named theorems in the active Lean
surface typecheck without local proof placeholders.  The analytic gap is not
hidden in Lean; it is explicitly named as a conditional hypothesis when it has
not been proved.

\section{Discussion}

The same-height framework contributes a structural diagnostic for RH proof
search.  It converts an off-line-zero problem into a degeneracy problem for the
height projection
\[
  s\mapsto \im(s)
\]
restricted to critical-strip zeros.  This view aligns with spectral intuition:
if zeta zeros were realized as a genuinely one-dimensional self-adjoint
spectrum, then same-height collisions off the critical line would be
geometrically unnatural.  But the framework does not assume such a spectral
model.

The most important limitation is also the main research target.  No currently
formalized theorem proves same-height repulsion for zeta zeros.  Statistical
spacing laws, finite computation, and known proportions of simple zeros provide
evidence for zero repulsion, but they do not supply the deterministic
fiberwise uniqueness statement needed here.  The failed convexity experiment is
therefore useful: it rules out an attractive but false sufficient condition and
pushes the program toward local horizontal argument-principle methods.

\section{Conclusion}

The paper presents a Lean-verified reduction framework for RH centered on
same-height zero fibers.  The central equivalence says that RH is exactly the
absence of same-height critical-strip zero collisions with different real
parts.  Additional verified criteria express the same obstruction through
fiber energy, fiber cardinality, vertical zero-free thresholds, and conditional
horizontal critical-point exclusions.  The development also formalizes a
\(3\)-\(4\)-\(1\) de la Vallee Poussin product bridge in the half-plane
\(\sigma>1\).

The next mathematical step is sharply defined: prove a local same-height
repulsion principle for \(\zetaR(x+i\gamma)\).  A plausible direction is a
horizontal argument-principle or componentwise Rolle theorem strengthened by
zeta-specific phase information.  Until that principle is proved, the work is a
verified reduction and proof-search architecture, not a proof of RH.

\section*{Acknowledgements}

The Lean development depends on mathlib and on its formalization of the
Riemann zeta function and the Riemann Hypothesis statement.

\bibliographystyle{alpha}
\bibliography{references}

\appendix

\section{Lean theorem map}

\begin{center}
\begin{tabular}{p{0.36\linewidth}p{0.54\linewidth}}
\hline
Mathematical statement & Lean theorem/module \\
\hline
RH iff same-height uniqueness &
\leanpath{riemannHypothesis_iff_forall_zeroFiberRealUnique_unconditional},
\leanpath{RiemannHypothesisReduction.lean} \\
RH iff no same-height collision &
\leanpath{riemannHypothesis_iff_noSameImaginaryPartCollision_unconditional},
\leanpath{RiemannHypothesisReduction.lean} \\
RH iff all fiber energies vanish &
\leanpath{riemannHypothesis_iff_forall_fiberEnergy_zero},
\leanpath{SpectralEnergyCriterion.lean} \\
Positive fiber energy detects RH failure &
\leanpath{not_riemannHypothesis_iff_exists_fiberEnergy_pos},
\leanpath{FiberEnergyGap.lean} \\
Fiber centroid equals \(1/2\) &
\leanpath{fiber_centroid_eq_half}, \leanpath{FiberCentroid.lean} \\
Fiber cardinality \(\leq 1\) implies RH &
\leanpath{riemannHypothesis_of_all_fiber_card_le_one},
\leanpath{FiberCentroid.lean} \\
Threshold zero-free equivalence &
\leanpath{riemannHypothesis_iff_zeroFreeRightOf_all_gt_half},
\leanpath{ZeroFreeEngine.lean} \\
\(3\)-\(4\)-\(1\) zeta product lower bound &
\leanpath{zeta_341_modulus}, \leanpath{Zeta341GlobalBridge.lean} \\
Convexity sufficient condition &
\leanpath{riemannHypothesis_of_horizontalStrictMidpointConvex},
\leanpath{HorizontalRepulsion.lean} \\
Componentwise critical exclusion sufficient condition &
\leanpath{riemannHypothesis_of_componentwiseCriticalExclusion},
\leanpath{HorizontalCritical.lean} \\
Spectral operator self-adjointness and zero-crossing exclusion &
\leanpath{SpectralPhaseBounds.lean.draft} \\
Topological flow line and saddle point exclusion &
\leanpath{TopologicalFlow.lean.draft} \\
Dynamical repulsion and outermost zero outward force &
\leanpath{DynamicalRepulsion.lean.draft} \\
\hline
\end{tabular}
\end{center}

\section{Submission note}

For peer review, the accurate submission category is a formalization or
proof-assistant-assisted reduction paper, not a claimed proof of the Riemann
Hypothesis.  A suitable arXiv classification would be \texttt{math.NT} with
cross-listing to \texttt{cs.LO} if the Lean development is included as a
reproducible artifact.

\end{document}
